% !TEX program = pdflatex
\documentclass[12pt,a4paper]{article}

\usepackage[T1]{fontenc}
\usepackage[margin=1in]{geometry}
\usepackage{xcolor}
\usepackage{hyperref}
\usepackage{setspace}
\usepackage{parskip}
\usepackage{enumitem}

\hypersetup{
  colorlinks=true,
  urlcolor=blue!70!black
}

\onehalfspacing
\pagestyle{empty}

\begin{document}

\begin{center}
{\Large\bfseries Personal Statement}\\[4pt]
{\large Odeyemi Olusegun Israel}\\[2pt]
{\normalsize UK Global Talent Visa --- Digital Technology (Exceptional Promise)}
\end{center}

\bigskip

I am an independent researcher and systems engineer based in Derby,
United Kingdom.  Over the past three years---without institutional
affiliation, funding, or supervision---I have designed and built a
blockchain-native compliance platform spanning 17~containerised
microservices and 18~smart contracts, trained machine-learning models on
2.64~million Ethereum transactions, produced three original research
papers establishing a unified mathematical theory of instability
detection across finance, energy infrastructure, and fluid dynamics,
filed a UK patent application covering 25~claims across 13~invention
groups, published on TechRxiv and Zenodo, and open-sourced developer
tools in TypeScript, Python, and Dart.  I am applying under the
Exceptional Promise route because my work demonstrates the capacity to
become a leading figure in the UK's applied mathematics, financial
technology, and machine learning sectors.

%% ── SECTION 1: ENGINEERING ─────────────────────────────────────────

\textbf{1.\enspace What I have built.}

My central engineering project is \textbf{AMTTP}---the Anti-Money
Laundering Transaction Trust Protocol.  It addresses a problem that
no existing commercial product has solved: enforcing regulatory compliance
\textit{before} a blockchain transaction reaches finality, rather than
generating alerts after the fact.  Chainalysis, Elliptic, and TRM~Labs
all operate post-settlement; AMTTP intervenes at the protocol level
within the approximately 12-second Ethereum block confirmation window.

\textit{Architecture.}\;
The system spans four layers.  At the outermost layer sit open-source
client SDKs in TypeScript (19+~service endpoints covering risk, KYC,
transactions, disputes, reputation, sanctions, GeoRisk, MEV~protection,
and bulk operations for up to 10,000~transactions) and Python
(17~matching services, full async API).  The compliance orchestration
layer comprises 7~parallel microservices---ML risk scoring, graph-based
relationship analysis, policy engine, sanctions screening, continuous
monitoring, geographic risk, and an explainability module---coordinated
through a central gateway with NGINX TLS termination, RS256 JWT
validation, per-key rate limiting, and Ed25519~oracle-signed responses
with 90-day rotation.  The ML pipeline (described below) supplies risk
scores to the orchestrator.  Finally, the smart contract layer enforces
on-chain compliance through 18~Solidity contracts deployed on Ethereum
Sepolia, including risk-tiered policy enforcement
(Approve / Review / Escrow / Block), Kleros-integrated decentralised
dispute resolution, LayerZero cross-chain policy synchronisation, Gnosis
Safe and Biconomy account-abstraction modules, and Compliance NFTs.  The
full infrastructure totals 17~containerised microservices---each with
health checks, circuit breaker patterns, and Prometheus / Grafana
monitoring---orchestrated through Docker Compose.

\textit{ML pipeline.}\;
The production fraud-detection model uses a Composite Teacher
architecture (40\% XGBoost, 30\% FATF pattern detectors, 30\% graph
structural properties) that generates pseudo-labels across
2.64~million Ethereum mainnet transactions sourced from Google BigQuery.
A stacked ensemble of XGBoost (500~estimators) and LightGBM
(500~estimators) with a logistic regression meta-learner achieves a
production ROC-AUC of \textbf{0.9879}, PR-AUC of \textbf{0.9715}, and
F1 of \textbf{0.9012} with 21~features and \textbf{zero
training-serving skew}.  An earlier v2 student pipeline integrating
BetaVAE (64 latent dimensions), GraphSAGE (3~layers), and GATv2
(4~attention heads) reaches ROC-AUC~0.9229.  Achieving zero skew was
itself a significant engineering challenge: v2 models trained on
71~features but served with only~5 (66~features hardcoded to zero);
diagnosing and correcting this lifted ROC-AUC from 0.64~to~0.99 in a
single iteration.

\textit{zkNAF.}\;
I also built \textbf{zkNAF}---a zero-knowledge Non-disclosing Anti-Fraud
framework using Groth16~ZK-SNARK proofs on the BN254~elliptic curve.
zkNAF provides four proof types: sanctions non-membership (Merkle
non-membership without revealing address or list), risk-range proofs
(ML~score within acceptable bounds without exposing the exact value),
KYC~verification (credential completeness without identity disclosure),
and combined transaction compliance.  This resolves the tension between
GDPR's data minimisation requirements and AML transparency obligations
at the cryptographic level, with dedicated on-chain verifier contracts
routed through a single \texttt{ZkNAFVerifierRouter}.

\textit{Consumer application.}\;
To make pre-settlement compliance accessible to end users, I built a
cross-platform Flutter consumer application targeting six platforms
(iOS, Android, Web, Windows, macOS, Linux) from a single Dart codebase.
The app includes a wallet dashboard, pre-transfer trust check with
traffic-light interstitial, transaction history, dispute filing, and
MetaMask wallet connection.  It communicates with the compliance backend
through a hybrid Flutter / Next.js architecture using a bidirectional
JavaScript Channel bridge.  A separate Next.js~15 ``War Room'' dashboard
serves institutional compliance analysts.

\textit{Security architecture.}\;
Security spans every layer: EIP-712 typed signing for request integrity,
server-enforced nonce + TTL (300\,s) replay protection via Redis,
HMAC-SHA256 webhook authentication, a 6-tier RBAC hierarchy, and a
UI Integrity Service that monitors DOM mutations and verifies per-component
hashes to prevent address-swap attacks of the kind used in the 2025
Bybit exploit.

The AMTTP architecture is described in my published paper:
\textit{Olusegun Odeyemi, ``AMTTP: A Four-Layer Architecture for
Deterministic Compliance Enforcement in Institutional DeFi,''
TechRxiv, February~27, 2026},
\href{https://doi.org/10.36227/techrxiv.177220113.33816607/v1}{DOI:\,10.36227/techrxiv.177220113.33816607/v1}.

%% ── SECTION 2: RESEARCH ────────────────────────────────────────────

\textbf{2.\enspace What I have discovered.}

Alongside the engineering work, I have pursued a research programme that
began with a practical question---why do fraud detectors miss what they
are designed to find?---and evolved into a unified mathematical theory of
instability in complex systems, with results spanning machine learning,
systemic financial risk, electrical grid failure, algorithmic stablecoin
collapse, and the regularity of the 3D Navier-Stokes equations.

\textit{Paper~1: Blind Spot Decomposition Theory (BSDT).}\;
My first paper provides a mathematical framework for characterising
missed fraud.  I decompose detection failures into four near-orthogonal
components---Camouflage~($\delta_C$), Feature Gap~($\delta_G$), Activity
Anomaly~($\delta_A$), and Temporal Novelty~($\delta_T$)---and prove that
these account for over 95\% of explainable variance in missed-fraud
indicators (component correlations $|\rho|<0.3$ for $>$80\% of pairs;
NMI~$<$~0.10 for $>$90\% of pairs).  The correction operators recover
recall by up to \textbf{84.4~percentage points} without retraining or
access to model internals.  Validated across 4~model architectures
(XGBoost, LightGBM, Random Forest, Logistic Regression), 7~datasets
from 6~domains (Bitcoin / Elliptic with 203,769~transactions, Ethereum
addresses, IEEE-CIS credit card, NASA Shuttle telemetry, mammography,
pen-digit recognition, e-commerce card-not-present), and over
36~model--dataset combinations, the best variant (F2\_QuadSurf) achieves
mean~F1 of~0.787 with 88.1\% false positive reduction.  The paper also
proves formal safety, boundedness, and monotonicity guarantees for the
correction formula and introduces three self-calibrating weight estimation
methods (mutual information, Fisher, Bayesian online).
Published on Zenodo (DOI:\,\href{https://doi.org/10.5281/zenodo.18870589}{10.5281/zenodo.18870589});
target journal: \textbf{IEEE~TKDE}.

\textit{Paper~2: Universal Deviation Law (UDL).}\;
My second paper tackles anomaly detection from first principles.
I introduce a multi-spectrum tensor framework with 14~composable
operators spanning five mathematical families---statistical, dynamical,
spectral, geometric, and exponential---each designed to capture a
distinct axis of deviation.  A $14\times14$ correlation audit confirms
that only 1~pair out of~91 is redundant (Wavelet$\leftrightarrow$Phase,
$\rho=0.92$): every other operator adds genuinely new information.
The best configuration achieves mean AUC of~\textbf{0.972} and 91\%
coverage across five benchmarks---categorically outperforming Isolation
Forest~(0.788), LOF~(0.632), and the current state-of-the-art
ECOD~(0.921)---with a single hyperparameter set and no per-dataset
tuning.  I prove five formal theorems providing local distinguishability
(via the Inverse Function Theorem), global sensitivity (uniform
bound~$\underline\sigma>0$), finite-sample concentration (Hoeffding),
SDP-optimal fusion weights, and a $(1-1/e)$-approximation guarantee for
greedy operator selection via monotone submodularity.  The paper also
introduces a GravityEngine N-body scoring module with a closed-form
erf attraction potential that achieves AUC~0.899 on mammography with
monotone energy descent.
Target venue: \textbf{NeurIPS~2026}.

\textit{Paper~3: Systemic Risk as Geometry (Adaptive Friction + Navier-Stokes).}\;
My third paper---and the centrepiece of the research programme---establishes
a geometric framework linking the BSDT deviation operators to the
Lyapunov energy of networked agent dynamics.  It proves three theorems
(equivalence of the detection functional and the physical potential;
a spectral phase transition at a critical manifold where constant
coupling becomes impossible; and a closed-form optimal damping policy
$\gamma^*(E) = E/(E+\theta)$) and validates them across four domains
using exclusively real-world data:

\begin{itemize}[nosep, leftmargin=2em]
  \item \textbf{Banking} (140~quarters of FDIC call-report data, 1990--2024):
    6-quarter lead before the 2008 GFC using a threshold trained
    exclusively on pre-2006 data; AUROC~$= 0.805$ at institution level
    for the 30 largest US banks.
    A 24-institution global panel spanning US, UK, EU, Asia, and Africa
    (using World Bank GFDD, ECB MIR, and BIS bilateral exposure data)
    achieves AUROC~$= 0.811$ with a 19-quarter GFC lead.  Threshold
    Granger causality passes on the G-SIB panel ($p = 0.043$).
    The framework also calibrates countercyclical capital buffers:
    the 2008~GFC consumption-equivalent welfare loss reached 19.57
    percentage points; the COVID peak CCyB calibration was 231~basis
    points.

  \item \textbf{Algorithmic stablecoins} (Terra/Luna UST, May 2022):
    detection at hour~23, one hour after the initial depeg event,
    with simulation-to-reality correlation $r = 0.996$.  Adaptive
    friction limits the depeg to 0.28\% (versus $>$99\% collapse
    without intervention).

  \item \textbf{Energy infrastructure} (Texas ERCOT, February 2021):
    detection at hour~24, a full $+$72~hours before rolling blackouts
    begin.  Adaptive friction reduces peak capacity loss by 71\%.
    The per-agent decomposition reveals that the cascade origin shifted
    from frozen wind turbines (the public narrative) to gas supply chain
    disruption (the structural vulnerability)---a finding invisible to
    event-study analysis.

  \item \textbf{Fluid dynamics} (3D incompressible Navier-Stokes):
    reinterpreting $\gamma^*$ as a state-dependent viscosity yields the
    \textbf{P/D halving principle}---adaptive viscosity halves the enstrophy
    production-to-dissipation ratio at every Reynolds number tested
    ($\mathrm{Re} \in [314, 6283]$), confirmed numerically on a
    pseudo-spectral solver for the Taylor-Green vortex.  This constitutes
    a new conditional regularity criterion connected to the Clay
    Millennium Prize Problem and identifies three strategies toward
    closing the gap to unconditional regularity (with explicit disclaimers
    that it does not solve the Clay problem).
\end{itemize}

\noindent
In all three non-financial domains, the dominant BSDT channel is
$\delta_C$ (camouflage)---the mechanism by which an unstable system
conceals its proximity to collapse---demonstrating that the mathematical
structure is genuinely universal, not fitted to any single domain.  The
unifying object is the blind-spot energy functional
$E_{\mathrm{BS}} = \delta_C^2 + \delta_G^2 + \delta_A^2 + \delta_T^2$,
which serves simultaneously as a fraud-detection diagnostic, a systemic
risk alarm, and a viscosity schedule for turbulent fluids.
Target journal: \textbf{SIAM Journal on Applied Mathematics}.

%% ── SECTION 3: IP ──────────────────────────────────────────────────

\textbf{3.\enspace Intellectual property.}

On 4~March~2026, I filed a UK patent application with the Intellectual
Property Office (Patents Act~1977) titled \textit{``Adaptive
Multi-Spectrum Anomaly Detection, Risk-Aware Transaction Processing,
and Cross-Domain Stability Control System.''} The application contains
25~claims organised into 13~groups covering: (A)~BSDT four-channel
decomposition and post-hoc correction, (B)~UDL full-rank detection
guarantees, (C)~GravityEngine N-body scoring, (D)~Dimension Magnifier
nonlinear transform, (E)~AMTTP risk-aware blockchain protocol with
4-tier escrow, (F)~zkNAF zero-knowledge compliance, (G)~cross-chain risk
propagation, (H)~UI integrity verification (anti-Bybit),
(I)~adaptive friction for systemic risk and CCyB calibration,
(J)~state-dependent viscosity for Navier-Stokes, (K)~compliance
orchestrator with parallel microservices, (L)~cross-domain unification
via the blind-spot energy functional, and (M)~decentralised dispute
resolution for ML-flagged transactions.

%% ── SECTION 4: OPEN-SOURCE ─────────────────────────────────────────

\textbf{4.\enspace Open-source contributions and tools.}

The full codebase---research code, pipeline implementations, client SDKs,
smart contracts, and replication packages---is publicly available on
GitHub at \texttt{github.com/segetii/fintech}.  Major open-source
components include:

\begin{itemize}[nosep, leftmargin=2em]
  \item \textbf{TypeScript client SDK} (\texttt{@amttp/client-sdk}):
    19+ service modules with EIP-712 typed signing, Flashbots MEV
    protection, HMAC-SHA256 webhook signing, and batch scoring for
    up to 10,000 transactions.
  \item \textbf{Python client SDK} (\texttt{amttp-sdk}):
    17~async service modules mirroring the TypeScript API, MIT licensed.
  \item \textbf{MFLS SDK} (\texttt{mfls}):
    complete implementation of the BSDT/MFLS systemic risk engine with
    BSDT operators, 5~scoring variants, GravityEngine, phase-transition
    detection, and a FastAPI~REST~API for systemic risk signal, blind-spot
    audit, CCyB calibration, and herding monitoring.
  \item \textbf{MetaMask Snap}: browser-based compliance verification
    for everyday users.
  \item \textbf{GravityEngine toolkit}: N-body simulation for networked
    agent dynamics with erf-attraction, log-repulsion potentials,
    spectral radius computation, and trajectory analysis.
  \item \textbf{3D Navier-Stokes solver}: pseudo-spectral solver with
    BSDT diagnostic suite and adaptive viscosity, including Taylor-Green
    vortex benchmarks across $\mathrm{Re}\in[314,6283]$.
  \item \textbf{Complete replication packages} for all empirical results
    in all three research papers.
\end{itemize}

%% ── SECTION 5: WHY UK ──────────────────────────────────────────────

\textbf{5.\enspace Why the UK, and why now.}

The UK is positioning itself as the global leader in responsible fintech
regulation.  The FCA's evolving framework for crypto-assets will create
substantial demand for compliance infrastructure that operates at the
protocol level.  AMTTP is, as far as I am aware from surveying the
commercial landscape, the only system that provides pre-settlement AML
enforcement on blockchain transactions.

The cross-domain reach of my research---from banking regulation to energy
grid resilience to pure mathematics---aligns with the UK's strategic
priorities in financial stability (Bank of England, PRA), infrastructure
resilience (Ofgem, National Grid ESO), and mathematical sciences (EPSRC,
the Heilbronn Institute).  The Navier-Stokes extension, in particular,
connects directly to the UK's world-leading position in applied
mathematics and fluid dynamics research.

My immediate plans are: (1)~submit the three papers to their target
journals (IEEE~TKDE, NeurIPS~2026, SIAM~J.~Applied Math);
(2)~prepare a short letter for \textit{Nature Communications}
summarising the cross-domain universality result for broader scientific
visibility; (3)~commercialise AMTTP as a compliance-as-a-service product
for UK-regulated financial institutions and virtual asset service
providers; and (4)~build a research-led company at the intersection of
machine learning, financial regulation, and mathematical stability
theory---headquartered in the UK.

%% ── CLOSING ────────────────────────────────────────────────────────

\medskip

I have done this work alone, without a university laboratory, a GPU
cluster, or a research budget.  In three years, I have gone from a
practical question about missed fraud to a unified mathematical framework
that detects collapse across banking systems, stablecoins, power grids,
and turbulent fluids---published two papers, filed a patent with
25~claims, deployed 18~smart contracts, trained models on 2.64~million
transactions to ROC-AUC~0.9879, and open-sourced the complete stack.
What I bring is the ability to move between rigorous mathematical theory
and production-grade systems, and the discipline to validate claims
before making them.  I believe I can contribute meaningfully to the UK's
digital economy, scientific reputation, and regulatory leadership, and
I am asking for the opportunity to do so.

\end{document}
